% ============================================================================
% WORKBOOK — Additional: Introduction to Square Roots
% State-specific (VA)
% Practice-focused reinforcement for the study guide topic
% ============================================================================
\section{Introduction to Square Roots}
\topicTitle{Introduction to Square Roots}

\begin{quickReview}{Perfect Squares and Square Roots}
A \textbf{perfect square} is a number you get by multiplying an integer by itself:

$1, 4, 9, 16, 25, 36, 49, 64, 81, 100, 121, 144, 169, 196, 225$

The \textbf{square root} undoes squaring: $\sqrt{25} = 5$ because $5 \times 5 = 25$.

\medskip
To \textbf{estimate} $\sqrt{n}$ for a non-perfect square, find the two perfect squares it falls between.

\textbf{Example:} $\sqrt{40}$: Since $36 < 40 < 49$, we know $6 < \sqrt{40} < 7$. Since $40$ is closer to $36$, $\sqrt{40} \approx 6.3$.
\end{quickReview}

% ── Warm-Up ──────────────────────────────────────────────────────────────────
\begin{practiceBox}[funGreen]{\faSun~Warm-Up}

\practiceHeader[funGreen]{Find the Square Root}

\begin{multicols}{2}
\prob $\sqrt{9} =$ \answerBlank[2cm]
\answerExplain{$3$}{$3 \times 3 = 9$, so $\sqrt{9} = 3$.}

\prob $\sqrt{49} =$ \answerBlank[2cm]
\answerExplain{$7$}{$7 \times 7 = 49$, so $\sqrt{49} = 7$.}

\prob $\sqrt{100} =$ \answerBlank[2cm]
\answerExplain{$10$}{$10 \times 10 = 100$, so $\sqrt{100} = 10$.}

\prob $\sqrt{36} =$ \answerBlank[2cm]
\answerExplain{$6$}{$6 \times 6 = 36$, so $\sqrt{36} = 6$.}

\prob $\sqrt{1} =$ \answerBlank[2cm]
\answerExplain{$1$}{$1 \times 1 = 1$, so $\sqrt{1} = 1$.}

\prob $\sqrt{64} =$ \answerBlank[2cm]
\answerExplain{$8$}{$8 \times 8 = 64$, so $\sqrt{64} = 8$.}
\end{multicols}
\end{practiceBox}

% ── Practice A: More Perfect Squares ─────────────────────────────────────────
\begin{practiceBox}{Perfect Square Roots}

\begin{multicols}{2}
\prob $\sqrt{121} =$ \answerBlank[2cm]
\answerExplain{$11$}{$11 \times 11 = 121$, so $\sqrt{121} = 11$.}

\prob $\sqrt{144} =$ \answerBlank[2cm]
\answerExplain{$12$}{$12 \times 12 = 144$, so $\sqrt{144} = 12$.}

\prob $\sqrt{169} =$ \answerBlank[2cm]
\answerExplain{$13$}{$13 \times 13 = 169$, so $\sqrt{169} = 13$.}

\prob $\sqrt{196} =$ \answerBlank[2cm]
\answerExplain{$14$}{$14 \times 14 = 196$, so $\sqrt{196} = 14$.}

\prob $\sqrt{225} =$ \answerBlank[2cm]
\answerExplain{$15$}{$15 \times 15 = 225$, so $\sqrt{225} = 15$.}

\prob $\sqrt{4} =$ \answerBlank[2cm]
\answerExplain{$2$}{$2 \times 2 = 4$, so $\sqrt{4} = 2$.}
\end{multicols}
\end{practiceBox}

% ── Practice B: Between Two Whole Numbers ────────────────────────────────────
\begin{practiceBox}[funTeal]{Between Which Two Whole Numbers?}

\practiceHeader[funTeal]{Find the two consecutive whole numbers that $\sqrt{n}$ falls between.}

\prob $\sqrt{10}$ is between \answerBlank[1cm] and \answerBlank[1cm].
\answerExplain{$3$ and $4$}{$9 < 10 < 16$, so $3 < \sqrt{10} < 4$.}

\prob $\sqrt{30}$ is between \answerBlank[1cm] and \answerBlank[1cm].
\answerExplain{$5$ and $6$}{$25 < 30 < 36$, so $5 < \sqrt{30} < 6$.}

\prob $\sqrt{50}$ is between \answerBlank[1cm] and \answerBlank[1cm].
\answerExplain{$7$ and $8$}{$49 < 50 < 64$, so $7 < \sqrt{50} < 8$.}

\prob $\sqrt{75}$ is between \answerBlank[1cm] and \answerBlank[1cm].
\answerExplain{$8$ and $9$}{$64 < 75 < 81$, so $8 < \sqrt{75} < 9$.}

\prob $\sqrt{110}$ is between \answerBlank[1cm] and \answerBlank[1cm].
\answerExplain{$10$ and $11$}{$100 < 110 < 121$, so $10 < \sqrt{110} < 11$.}
\end{practiceBox}

% ── Practice C: Estimating ───────────────────────────────────────────────────
\begin{practiceBox}[funOrange]{Estimate to One Decimal Place}

\practiceHeader[funOrange]{Estimate each square root to one decimal place.}

\prob $\sqrt{20} \approx$ \answerBlank[2cm]
\answerExplain{$\approx 4.5$}{$16 < 20 < 25$. $20$ is about halfway, so $\sqrt{20} \approx 4.5$.}

\prob $\sqrt{5} \approx$ \answerBlank[2cm]
\answerExplain{$\approx 2.2$}{$4 < 5 < 9$. $5$ is very close to $4$, so $\sqrt{5} \approx 2.2$.}

\prob $\sqrt{45} \approx$ \answerBlank[2cm]
\answerExplain{$\approx 6.7$}{$36 < 45 < 49$. $45$ is close to $49$, so $\sqrt{45} \approx 6.7$.}

\prob $\sqrt{90} \approx$ \answerBlank[2cm]
\answerExplain{$\approx 9.5$}{$81 < 90 < 100$. $90$ is about halfway, so $\sqrt{90} \approx 9.5$.}

\prob $\sqrt{150} \approx$ \answerBlank[2cm]
\answerExplain{$\approx 12.2$}{$144 < 150 < 169$. $150$ is close to $144$, so $\sqrt{150} \approx 12.2$.}
\end{practiceBox}

% ── True or False ────────────────────────────────────────────────────────────
\begin{practiceBox}[funPurple]{True or False?}

\trueOrFalse{$\sqrt{16} = 8$}
\answerExplain{False}{$\sqrt{16} = 4$ because $4 \times 4 = 16$. The number $8$ is not the square root of $16$.}

\trueOrFalse{Every perfect square has a whole number as its square root.}
\answerExplain{True}{By definition, a perfect square is the product of a whole number multiplied by itself, so its square root is that whole number.}

\trueOrFalse{$\sqrt{2}$ is between $1$ and $2$.}
\answerExplain{True}{$1 < 2 < 4$, so $1 < \sqrt{2} < 2$.}
\end{practiceBox}

% ── Word Problems ────────────────────────────────────────────────────────────
\begin{practiceBox}[funBlue]{\faGlobe~Word Problems}

\wordProblem{A square garden has an area of $81$ ft$^2$. What is the side length?}{ft}
\answerExplain{$9$ ft}{Side $= \sqrt{81} = 9$ ft because $9 \times 9 = 81$.}

\wordProblem{A square tile has an area of $196$ cm$^2$. What is the side length?}{cm}
\answerExplain{$14$ cm}{Side $= \sqrt{196} = 14$ cm because $14 \times 14 = 196$.}

\wordProblem{A square room has an area of $130$ ft$^2$. Estimate the side length to one decimal place.}{ft}
\answerExplain{$\approx 11.4$ ft}{$121 < 130 < 144$, so $11 < \sqrt{130} < 12$. Since $130$ is closer to $121$, $\sqrt{130} \approx 11.4$ ft.}
\end{practiceBox}

% ── Challenge ────────────────────────────────────────────────────────────────
\begin{challengeBox}
\prob Is $\sqrt{200}$ closer to $14$ or $15$? Explain your reasoning. \answerBlank[3cm]
\answerExplain{Closer to $14$}{$14^2 = 196$ and $15^2 = 225$. Since $200$ is much closer to $196$ than to $225$, $\sqrt{200}$ is closer to $14$ (about $14.1$).}
\end{challengeBox}

\encouragement{You are mastering square roots! Knowing your perfect squares makes estimating a breeze!}
